%%%%% templates/article/basic.tex
%%%%% preambles/package-importing.tex
% Paper margins
\usepackage{geometry}

% AMS package group
\usepackage{amsmath,amsfonts,amssymb,amsthm}

% Math extended derivative package
\usepackage{derivative}

% figure inserting
\usepackage{graphicx}
\usepackage{float}


% Colors
\usepackage[dvipsnames, svgnames]{xcolor}

% SyntaxHighlighting
\usepackage{listings}

% Paragraph indent
\usepackage{indentfirst}

% Hyperlinks
\usepackage[bookmarks=true,unicode]{hyperref}
\usepackage{bookmark}
\usepackage{hyperxmp}

% Table stuffs
% Used for table cell settings
\usepackage{array, multirow, booktabs, longtable, makecell}

% inserting stand alone files inside.
\usepackage{standalone}

% TIKZ-- careful, this would require some libraries to become funtional
\usepackage{tikz}

\usepackage{pgfplots}
\pgfplotsset{compat=1.18}

\usepackage{mathrsfs}
% <tikz libraries>
\usetikzlibrary{calc, shadows, positioning, intersections}
\usetikzlibrary{arrows.meta, automata, arrows}
\usetikzlibrary{angles, quotes}
\usetikzlibrary{matrix}
% </tikz libraries>

% Multiline commenting
\usepackage{comment}

% extra plainer symbols that can be used to fake math content inside pdf bookmarks
\usepackage{textcomp}

% To include .pdf_tex, which is essentially made by inkscape from .svg files, to the .tex document
\usepackage{import,xifthen,pdfpages}

% enumitem for enumeration options (e.g. [label={\alph*}] to make it alphabetical)
\usepackage{enumitem}

% SI units extention. Using it by: \unit{}
\usepackage{siunitx}

% Extensive font specifications
\usepackage{fontspec}

\usepackage{xparse}

% font-sizing
\usepackage[fontsize=11pt]{fontsize}
\usepackage{anyfontsize}

%\usepackage[T1]{inputenc}


%==================================================%

%%%%% preambles/cross-ref.tex
\hypersetup{
  hidelinks,
%<pdf>
  bookmarksnumbered,
  pdfmenubar=true,
%</pdf>
%<colors>
  colorlinks=true,
  citecolor=OliveGreen,
  filecolor=cyan,
  linkcolor=NavyBlue,
  urlcolor=RedViolet,
%</colors>
  linktoc=all,
  pdfstartview=Fit,
  breaklinks=true
}


%==================================================%

%%%%% preambles/syntax-highlighting.tex
% requires: lstlistings, xcolor

%%%%% <colors>
\definecolor{codegreen}{rgb}{0,0.6,0}
\definecolor{codegray}{rgb}{0.5,0.5,0.5}
\definecolor{codepurple}{rgb}{0.58,0,0.82}
\definecolor{backcolour}{rgb}{0.95,0.95,0.92}
\definecolor{lightgray}{rgb}{.9,.9,.9}
\definecolor{darkgray}{rgb}{.4,.4,.4}
\definecolor{purple}{rgb}{0.65, 0.12, 0.82}
%%%%% </colors>

%%%%% <styles>
% Copied from: https://www.overleaf.com/learn/latex/Code_listing

\lstdefinestyle{overleafSampleCode}{
  backgroundcolor=\color{backcolour},  % code background
  commentstyle=\color{codegreen}, % comment
  keywordstyle=\color{magenta}, % language keywords
  numberstyle=\tiny\color{codegray}, % line-number
  stringstyle=\color{codepurple},
  basicstyle=\ttfamily\footnotesize,
  %%%%% <insert language="JavaScript">
  ndkeywordstyle=\color{darkgray}\bfseries,
  identifierstyle=\color{purple},
  extendedchars=true,
  %%%%% </insert>
  breakatwhitespace=false,         
  breaklines=true,                 
  captionpos=b,                    
  keepspaces=true,                 
  numbers=left,     % where to put the line-numbers; possible values are (none, left, right)               
  numbersep=5pt,   % how far the line-numbers are from the code               
  showspaces=false,                
  showstringspaces=false,
  showtabs=false,   % sets default tabsize to 2 spaces               
  tabsize=2,
  frame=single,	          % adds a frame around the code
}

% Copied from: https://tex.stackexchange.com/questions/89574/language-option-supported-in-listings/684030#684030?newreg=b1c105c269e54815b08ebc91f3fc811e
% Javascript extra:
\lstdefinelanguage{JavaScript}{
  keywords={abstract, any, as, boolean, break, case, catch, class, console, const, continue, debugger, declare, default, delete, do, else, enum, export, extends, false, finally, for, from, function, get, if, implements, import, in, infer, instanceof, interface, keyof, let, module, namespace, never, new, null, number, object, package, private, protected, public, readonly, require, return, set, static, string, super, switch, symbol, this, throw, true, try, type, typeof, undefined, unique, unknown, var, void, while, with, yield},
  morecomment=[l]{//},
  morecomment=[s]{/*}{*/},
  morestring=[b]',
  morestring=[b]",
  ndkeywords={class, export, boolean, throw, implements, import, this},
  sensitive=true
}

%%%%% </styles>
\lstset{style=overleafSampleCode}

%==================================================%


%%%%% preambles/macros/math-commands.tex
\newcommand{\vecbf}[1]{\ensuremath{\vec{\mathbf{#1}}}}
\newcommand{\Mod}[1]{\ (\mathrm{mod}\ #1)}

\newcommand{\tefrac}[2]{\frac{\text{#1}}{\text{#2}}}
\newcommand*\keystroke[1]{ % #REF -  https://tex.stackexchange.com/questions/5226/keyboard-font-for-latex
  \tikz[baseline=(key.base)]
    \node[%
      draw,
      fill=white,
      drop shadow={shadow xshift=0.25ex,shadow yshift=-0.25ex,fill=black,opacity=0.75},
      rectangle,
      rounded corners=2pt,
      inner sep=1pt,
      line width=0.5pt,
      font=\scriptsize\ttfamily
    ](key) {#1\strut}
  ;
}


\newenvironment{alieqn}{%
  \begin{equation}%
    \begin{aligned}%
}{%
    \end{aligned}%
  \end{equation}%
}

% Math shortcuts
\newcommand{\so}{\therefore}
\newcommand{\due}{\because}
\newcommand{\fa}{\forall}

% math text command
\newcommand{\tx}[1]{\text{#1}}
% parenthesis automatic L&R
\newcommand{\of}[1]{\left(#1\right)}
% square brackets automatic L&R
\newcommand{\off}[1]{\left[#1\right]}
% curly braces automatic L&R
\newcommand{\offf}[1]{\left\{#1\right\}}
% 
\newcommand{\at}[1]{\left.#1\right|}


%==================================================%

%%%%% preambles/macros/short-svg-input.tex
% \incfig is used for .pdf_tex files generated from .svg files, which makes .svg embedable into LaTeX

\NewDocumentCommand{\incfig}{%
  O{1} O{./figures/} m%
}{%
  \def\svgwidth{#1\columnwidth}%
  \import{#2}{#3.pdf_tex}%
}


%\newcommand{\incfig}[2][1]{%
%  \def\svgwidth{#1\columnwidth}%
%  \import{./figures/}{#2.pdf_tex}%
%}



%==================================================%

%%%%% preambles/macros/short-table-commands.tex
\usepackage{diagbox}
\usepackage{slashbox}
\newcommand*{\dl}{\\ \hline}
\newcolumntype{P}[1]{>{\centering\arraybackslash}p{#1}}
\newcolumntype{M}[1]{>{\arraybackslash}m{#1}}
\newcolumntype{C}[1]{>{\centering\arraybackslash}m{#1}}
\newcolumntype{L}[1]{>{\raggedright\arraybackslash}m{#1}}

\newcommand{\strike}{\backslashbox{\phantom{.}}{\phantom{.}}} % from: https://www.overleaf.com/latex/examples/table-cells-with-diagonal-lines/zrnpkbmzvtkr

% for longtable environments::
% \caption{}\label{} \dl
% headers~~~ \dl
% \endfirsthead
% \multicolumn{}{c}{\sffamily\large 续表}
% \endhead
% \bottomrule  
% \endfoot


%==================================================%

%%%%% preambles/macros/alignment.tex
\newenvironment{centereq}[1][0.8]
{% Source: https://tex.stackexchange.com/questions/576930/center-with-equal-amount-of-text-per-line
\trivlist
\centering
\setlength{\leftskip}{\dimexpr(\columnwidth-#1\columnwidth)/2\relax plus 1em}%
\setlength{\rightskip}{\leftskip}%
\item\relax
}
{\endtrivlist}

% paragraph indentation
\setlength{\parindent}{1.5em}

% paragraph separation
\setlength{\parskip}{\baselineskip}

%\newcommand{\q}[1]{{\par\vspace*{0.4cm}\noindent\large\bfseries #1 \par\vspace*{0.4cm}}}
%\newcommand{\sq}[1]{{\par\vspace*{0.4cm}\noindent\large\itshape #1 \par\vspace*{0.4cm}}}

\newcommand{\appendixpagenumbering}{
  \break
  \pagenumbering{arabic}
  \renewcommand{\thepage}{\thesection-\arabic{page}}
}

%==================================================%


%%%%% preambles/fonts/simplified-chinese.tex
\usepackage[fontset=none, scheme=plain]{ctex}
%%%%% <main font>
% mainfont：{Founder:方正新书宋_GBK	Regular}

% boldfont: {Founder:方正粗宋简繁_GBK}

% italicfont: {Founder:方正汉文正楷_GBK}



\setCJKmainfont[Path={/mnt/c/Windows/Fonts/},BoldFont={FZCSJW.TTF}, ItalicFont={FZKTK}]{FZXSSK}
%%%%% </main font>

%%%%% <sans serif font>


% sansfont: 方正黑体_GBK
% boldfont: 方正粗黑宋_GBK

\setCJKsansfont[Path={/mnt/c/Windows/Fonts/},BoldFont={FZCHSJW.TTF}]{FZHTK.TTF}

%%%%% </sans serif font>

% <font-families>
% 仿宋字体：{Founder:方正仿宋_GBK	Regular}
%   标识符：FZFSK--GBK1-0
%   文件：FZFSK.TTF
%\setCJKmonofont{FZFSK--GBK1-0}
%\setCJKfamilyfont{zhfangsong}{FZFSJW--GB1-0}
%\newcommand*{\fangsong}{\CJKfamily{zhfangsong}}




% </font-families>
\ctexset{space=auto} % Automatic chinese & english word spacing


%==================================================%


%%%%% preambles/class-specific/article.tex
% Refer to the last page. This one also stuck with the beamer class
%\usepackage{lastpage}

% Changing the styles of sections.
\usepackage{titlesec}

% Header&Footer
\usepackage{fancyhdr}

% This one would stop latex beamer from running
\usepackage{caption2}

\geometry{left=1in,right=1in,top=.8in,bottom=.8in}

%==================================================%


\standaloneconfig{mode=build}
\linespread{1.5}

%==================================================%

